% ScholarWeave LaTeX export — "document" doc type.
%
% This is NOT a full pandoc --template (see book.tex's own note for why) — it
% is included into the document PREAMBLE via pandoc's --include-in-header,
% alongside pandoc's own default LaTeX template (which already renders the
% TOC, citeproc bibliography, and font-package setup — no need to reproduce
% any of that here).
%
% Every line below is ordinary, editable LaTeX preamble content — copy this
% file (e.g. into your own Export Templates folder) and change, add, or
% remove anything: a journal logo (see \swTitleLogo below), 2-column layout
% (see SWEXTRACLASSOPTIONS below), different heading styles, extra packages,
% etc. Python only ever supplies the note's CONTENT (title/subtitle/author/
% date/abstract/body/citations) — never how it looks.
%
% Available substitutions, filled in by DocumentCompiler.py before this file
% reaches pandoc (NOT pandoc $variable$ syntax — see the note below):
%   SWTOKAUTHOR       — the note's author
%   SWTOKSHORTTITLE   — the note's short title (falls back to the full title)
%   SWTOKTITLE        — the note's title (unused by default; available if
%                        you want it somewhere other than the title block,
%                        e.g. in a running header instead of the short title)
% Available hooks (LaTeX commands DocumentCompiler.py invokes when building
% the title block/page — redefine them below to customize; left empty they
% do nothing):
%   \swTitleLogo      — invoked once, centered, just above the title
% Available marker (a plain LaTeX comment DocumentCompiler.py itself reads —
% inert to LaTeX, meaningful only to the Python side):
%   % SWEXTRACLASSOPTIONS: <comma-separated \documentclass options>
%                        — e.g. "twocolumn" or "12pt". "twoside" is always
%                        added too (required by the running-header
%                        convention below) — you don't need to repeat it.
% Available substitution tokens driven by export checkboxes, not user-edited
% (listed for completeness — moving/removing them changes what those
% checkboxes do):
%   SWTOKFIGURECOUNTER   — "Restart footnote and figure numbering per
%                        chapter": \counterwithin{figure}{section} (document/
%                        article's top-level heading is \section, not
%                        \chapter) when on, empty when off. Footnotes are
%                        deliberately NOT part of this: article-class
%                        footnotes are naturally continuous with no native
%                        per-section reset to hook into (unlike book, where
%                        \chapter resets footnotes to plain "1, 2, 3..."
%                        automatically — see book.tex), and continuous is
%                        the desired behaviour here regardless of the
%                        checkbox, so there's nothing to substitute.
%   SWTOKNEWPAGEHEADING  — "Top-level headings start on a new page": \newpage
%                        when on, empty when off. Sits in \titleformat's
%                        before-code slot for \section only (Heading 2 never
%                        forces a page break, matching DOCX/ODT).
%
% SWTOKAUTHOR / SWTOKSHORTTITLE / SWTOKTITLE are replaced with the document's
% real (LaTeX-escaped) values before this file is handed to pandoc — NOT
% pandoc $variable$ syntax, which LaTeX would read as inline math ($..$)
% since --include-in-header content is raw LaTeX, not run through pandoc's
% templating engine.
%
% Title BLOCK, not title page: matches document.docx/document.odt, whose
% title/subtitle/author/date sit as ordinary paragraphs at the top of page 1
% with no page break before the body — DocumentCompiler.py builds this as
% body content (see export_latex's use of is_book to choose title-block vs.
% title-page). "article" gets the same treatment; only "book" gets a
% dedicated title page, matching book.docx's own section-break-separated
% cover page.
%
% Running header: Author (even pages) / Short Title (odd pages), matching
% the DOCX/ODT "document"/"article" templates' own even/odd header
% convention — requires a two-sided layout (set from Python via
% --metadata classoption=twoside). No header on page 1: LaTeX's kernel
% already puts a document's first page on the `plain` page style regardless
% of \pagestyle, so this falls out for free.

% SWEXTRACLASSOPTIONS:

% ── Customize this template ────────────────────────────────────────────
% Copy this file (e.g. into your own Export Templates folder) and edit any
% of the four lines below to get your own fonts/margins/link color/spacing.
\setmainfont{Noto Sans}[BoldFont={Noto Sans},BoldFeatures={FakeBold=3.5}]  % body text font (must be installed; matches document.docx/.odt) — FakeBold: Noto Sans ships as ONE variable-weight file with no separately-loadable true bold instance XeTeX can select (confirmed: even explicitly requesting its Bold named instance silently produces the same regular-weight glyphs, with no warning), so \textbf/\bfseries need a synthetic (stroke-emboldened) bold instead of a real one
\usepackage[letterpaper,margin=1.25in]{geometry}  % page size + margins (try a4paper, or a per-side margin like "top=1in,bottom=1in,left=1.25in,right=1in")
\definecolor{swlinkcolor}{HTML}{156082}        % link/citation color, as a hex RRGGBB (no '#') — matches document.docx/.odt's Hyperlink style
\usepackage{setspace}\setstretch{1.15}         % paragraph line spacing (1 = single; matches document.docx/.odt's 1.15)
\raggedbottom                                  % let pages end wherever content naturally does, instead of
                                                % stretching heading/list/paragraph spacing to reach an even
                                                % bottom margin on every page (that stretch is what makes the
                                                % same heading's gap look bigger on one page than another with
                                                % no visible cause). Switch to \flushbottom for even bottom
                                                % margins instead, at the cost of that inconsistent stretch.
% ────────────────────────────────────────────────────────────────────────

% Title-block logo hook — empty by default. Redefine to add a journal/
% institution logo above the title, e.g.:
%   \renewcommand{\swTitleLogo}{\includegraphics[width=3cm]{/path/to/logo.png}\\[1em]}
\newcommand{\swTitleLogo}{}

% Arabic poetry font — "Noto Serif" above only covers Latin/Cyrillic/Greek
% (confirmed by inspecting its own cmap: zero Arabic codepoints), so Arabic
% text needs a font that actually has Arabic glyphs. The [Script=Arabic,RTL]
% fontspec feature key pair is what actually makes this correct — tried
% three approaches, in order:
%   1. A plain \newfontfamily with no direction info: letters shape/join
%      correctly (XeTeX's HarfBuzz-based shaper recognises the Arabic script
%      run on its own), but the run still lays out left-to-right — shaping
%      alone doesn't reorder. Looked right in an isolated render, but real
%      documents showed the wrong order in both Preview.app and Adobe Reader.
%   2. \TeXXeTstate=1 (XeTeX's native bidi engine) + wrapping each verse line
%      in \beginR...\endR (sw-poetry.lua): fixed the ORDER, but broke the
%      JOINING — letters came out disjointed/isolated-form (confirmed on a
%      real export). TeX--XeT reorders character nodes below the level where
%      Arabic contextual shaping happens, so the two don't compose.
%   3. [Script=Arabic,RTL] below: fontspec's own documented mechanism for a
%      short RTL insertion in an LTR document — applies shaping and ordering
%      together, correctly, needs no \beginR/\endR or \TeXXeTstate at all,
%      and (stress-tested against 60 footnotes) doesn't have polyglossia's
%      \begin{Arabic}...\end{Arabic} problem ("TeX capacity exceeded" on a
%      real multi-footnote document, confirmed earlier) since it never
%      engages polyglossia's document-wide language-switching machinery.
% Change the font name here to use a different Arabic typeface (e.g. "Noto
% Sans Arabic").
\newfontfamily\arabicfont[Script=Arabic,RTL]{Noto Naskh Arabic}

% Bullet-list levels alternate bullet/dash instead of LaTeX's own default
% bullet/dash/asterisk/centered-dot cascade.
\usepackage{enumitem}
\setlist[itemize,1]{label=\textbullet}
\setlist[itemize,2]{label=\textendash}
\setlist[itemize,3]{label=\textbullet}
\setlist[itemize,4]{label=\textendash}

% Heading 1 (\section) / Heading 2 (\subsection) styles — edit these to
% change appearance (bold/size/spacing/etc.); deeper heading levels use
% LaTeX's own defaults. No number prefix by default (\thesection etc. exist
% if you want to add one back — see titlesec's own documentation).
%
% \titlespacing's before/after values are FIXED lengths (no "plus/minus"
% stretch component) deliberately: the class's own default heading spacing
% is stretchy glue, which LaTeX's page-breaking algorithm will expand to
% help a page reach the bottom margin evenly — so the gap above a heading
% can end up looking much larger on one page than another with no visible
% cause. Fixed lengths make heading spacing the same everywhere, at the
% cost of pages no longer stretching to a perfectly even bottom margin.
\usepackage{titlesec}
\titleformat{\section}{\normalfont\Large\bfseries}{}{0pt}{SWTOKNEWPAGEHEADING}
\titlespacing*{\section}{0pt}{3.5ex}{2.3ex}
\titleformat{\subsection}{\normalfont\large\bfseries}{}{0pt}{}
\titlespacing*{\subsection}{0pt}{3.25ex}{1.5ex}
% Headings show no number (above), and the TOC entries below match — but
% \thesection still increments normally (--metadata numbersections=true,
% set in DocumentCompiler.py): SWTOKFIGURECOUNTER below needs a real,
% incrementing section counter to scope figures by, which secnumdepth's
% usual "no visible number" approach silently prevents (it stops the
% counter, not just the display).
\usepackage{titletoc}
\titlecontents{section}[0pt]{}{}{}{\hfill\thecontentspage}
\titlecontents{subsection}[1.5em]{}{}{}{\hfill\thecontentspage}
SWTOKFIGURECOUNTER

\usepackage{fancyhdr}
\pagestyle{fancy}
\fancyhf{}
\fancyhead[LE]{SWTOKAUTHOR}
\fancyhead[RO]{SWTOKSHORTTITLE}
\fancyfoot[C]{\thepage}
\renewcommand{\headrulewidth}{0.4pt}
% Page 1 (the title block) uses LaTeX's "plain" pagestyle — redefine it to
% keep the centered page-number footer but drop the running header, since
% \_latex_title_block sets \thispagestyle{plain} explicitly (article class
% has no automatic page-1 pagestyle override the way \maketitle would give
% it, since we never call \maketitle — see the title-block note above).
\fancypagestyle{plain}{%
  \fancyhf{}
  \fancyfoot[C]{\thepage}
  \renewcommand{\headrulewidth}{0pt}
}
% Not setting \hypersetup{pdftitle=...} here deliberately: --include-in-header
% content lands BEFORE pandoc's own hyperref setup in its default template, so
% \hypersetup isn't defined yet at this point — the PDF's internal
% pdftitle/pdfauthor metadata is the one thing not replicated for LaTeX
% output (a cosmetic gap, not a content one). colorlinks/linkcolor above
% DOES reach pandoc's own hypersetup call, since that reads --metadata
% variables Python passes separately, resolved against swlinkcolor once
% xcolor's \definecolor above has already run.
